% Motivation letter — Germany, TU Berlin / Studienkolleg pathway
% Adapt for: TU-Studienkolleg at TU Berlin, FU Berlin, or another T-Kurs.

% Shared LaTeX preamble for all motivation letters.
% Don't compile this file directly --- it's imported by each per-university .tex.
\documentclass[11pt,a4paper]{letter}

\usepackage[a4paper,margin=2.2cm]{geometry}
\usepackage[T1]{fontenc}
\usepackage[utf8]{inputenc}
\usepackage[sfdefault]{roboto}
\usepackage[hidelinks]{hyperref}
\usepackage{xcolor}
\usepackage{microtype}
\usepackage{parskip}

\definecolor{accent}{HTML}{0e7490}

\hypersetup{
  colorlinks=true,
  urlcolor=accent,
}

% Legal name used here — required on university applications.
% Professional name "Ahmad Bagheri" is used on LinkedIn / resume / OSS.
\name{Ahmad Baghereslami}
\address{
  contact@ahmadcodes.com \\
  +37494308785 \\
  ahmadcodes.com \\
  Armenia
}
\signature{Ahmad Baghereslami}

\input{glyphtounicode}
\pdfgentounicode=1


\begin{document}

\begin{letter}{
  Akademisches Auslandsamt\\
  Technische Universität Berlin\\
  Straße des 17.\ Juni 135\\
  10623 Berlin, Germany
}

\opening{Dear Selection Committee,}

I am applying for the T-Kurs at the TU-Studienkolleg, with the intent to enrol in the Bachelor's in Informatics (Informatik) at the Technische Universität Berlin in winter semester 2028 after passing the Feststellungsprüfung.

I am a senior software engineer with seven years of professional experience. Currently I am the sole engineer building Woody Portal at Barriertek, a fire-retardant lumber producer in Toronto, replacing phone-tag and whiteboard production scheduling with a unified portal: NestJS API, React + TanStack Query frontend, bi-directional NetSuite integration via SuiteScript RESTlets. PostgreSQL on Linux, Coolify-managed deployment. The portal is live with internal users; full external rollout is in October 2026. I have open-sourced the integration client at \href{https://github.com/Ahmad-b1995/netsuite-restlet-typescript}{github.com/Ahmad-b1995/netsuite-restlet-typescript}.

Before Barriertek I led a three-developer frontend team at DexTrading, a cryptocurrency analytics platform, working on Next.js with TanStack Query and shipping a NestJS + GraphQL content layer. Earlier I worked at Chargoon, an Iranian ERP vendor, migrating legacy modules to React and .NET as part of a 15-engineer team. Before that I shipped a PWA video streaming product at Hamisheh using Angular and HLS.

I attended the Humanities track at high school in Iran and earned my Pre-University Certificate in 2014 without taking the Konkur entrance examination. I began working as a programmer instead. Seven years later I have the practical skills a Bachelor's programme would teach, but I lack the structured theoretical grounding --- in mathematics, algorithms, and formal computer science --- that I now want to complete properly.

I am applying through the Studienkolleg pathway specifically because I understand the German system's expectation of subject-matter preparation before university admission, and the T-Kurs is the most honest way for someone with my profile (humanities-track high school, professional career in computer science) to bridge into a German Bachelor's. I view the year of preparatory study not as a delay but as the foundation I never built. German language acquisition is part of that --- I am currently at A1 and committed to reaching the C1 required for the Feststellungsprüfung.

TU Berlin specifically attracts me for three reasons. First, the Institut für Softwaretechnik und Theoretische Informatik publishes work that intersects directly with my professional domain (model-driven engineering for enterprise systems). Second, Berlin's professional ecosystem --- a large Iranian community, a strong open-source culture, and a software engineering job market that values practical experience --- makes it the right city for me and my partner to relocate to. Third, TU's emphasis on combining theoretical depth with applied research matches how I learn best.

During the Studienkolleg and the Bachelor's that follows, I plan to continue open-source work in ERP developer tooling and selective contract engagements in my specialty within the working-hour limits of the student visa. After graduation, I aim to build a consulting practice in ERP integration for German manufacturing and distribution businesses going through digital transformation --- a market that closely matches the work I already do, with the added grounding to do it more rigorously.

I would be grateful for the opportunity to study at the TU-Studienkolleg and eventually at TU Berlin.

\closing{Mit freundlichen Grüßen,}

\end{letter}

\end{document}
